% !PRINTED 191
\weberpage{191}

If we divide these formulae by $\vartheta_{01}(u)^{n^2}$, the
right-hand sides can be represented as integral rational functions
of the elliptic functions $\mathrm{sn}\,v$, $\mathrm{cn}\,v$,
$\mathrm{dn}\,v$ (§ 42). If, then, we set

\weq{3}{v = \pi\vartheta_{00}^2u, \quad \mathrm{sn}\,v=x, \quad \mathrm{cn}\,v=y, \quad \mathrm{dn}\,v=z}

we can write the formulae (1), (2), on determining a constant
factor suitably, as follows:

I. for even $n$:

\[
\begin{array}{ll}
\dfrac{\vartheta_{01}^{n^2}\vartheta_{00}}{\vartheta_{10}}\,\dfrac{\vartheta_{11}(nu)}{\vartheta_{01}(u)^{n^2}} = xyz\,A(x^2), & A(0)=n,\\[3mm]
\dfrac{\vartheta_{01}^{n^2}}{\vartheta_{10}}\,\dfrac{\vartheta_{10}(nu)}{\vartheta_{01}(u)^{n^2}} = B(x^2), & B(0)=1,\\[3mm]
\dfrac{\vartheta_{01}^{n^2}}{\vartheta_{00}}\,\dfrac{\vartheta_{00}(nu)}{\vartheta_{01}(u)^{n^2}} = C(x^2), & C(0)=1,\\[3mm]
\vartheta_{01}^{n^2-1}\,\dfrac{\vartheta_{01}(nu)}{\vartheta_{01}(u)^{n^2}} = D(x^2), & D(0)=1.
\end{array}
\]

II. for odd $n$:

\[
\begin{array}{ll}
\dfrac{\vartheta_{00}\vartheta_{01}^{n^2-1}}{\vartheta_{10}}\,\dfrac{\vartheta_{11}(nu)}{\vartheta_{01}(u)^{n^2}} = x\,A(x^2), & A(0)=n,\\[3mm]
\dfrac{\vartheta_{01}^{n^2}}{\vartheta_{10}}\,\dfrac{\vartheta_{10}(nu)}{\vartheta_{01}(u)^{n^2}} = y\,B(x^2), & B(0)=1,\\[3mm]
\dfrac{\vartheta_{01}^{n^2}}{\vartheta_{00}}\,\dfrac{\vartheta_{00}(nu)}{\vartheta_{01}(u)^{n^2}} = z\,C(x^2), & C(0)=1,\\[3mm]
\vartheta_{01}^{n^2-1}\,\dfrac{\vartheta_{01}(nu)}{\vartheta_{01}(u)^{n^2}} = D(x^2), & D(0)=1,
\end{array}
\]

where $A$, $B$, $C$, $D$ are integral rational functions of $x^2$,
whose degrees follow from the formulae (1), (2) in this way:

% [sic] In the odd-n row the German source prints (n^2-2)/2 for the degree of D, which is nonintegral for odd n; the defining formula implies (n^2-1)/2.
\[
\begin{array}{c@{\quad}cccc@{\qquad}l}
 & A(x^2), & B(x^2), & C(x^2), & D(x^2), & \\[3mm]
\text{Degree:} & \dfrac{n^2}{2}-2, & \dfrac{n^2}{2}, & \dfrac{n^2}{2}, & \dfrac{n^2}{2} & n \text{ even,}\\[3mm]
 & \dfrac{n^2-1}{2}, & \dfrac{n^2-1}{2}, & \dfrac{n^2-1}{2}, & \dfrac{n^2-2}{2} & n \text{ odd.}
\end{array}
\]

From the definitions I., II. one obtains immediately the values of
$A(0)$, $B(0)$, $C(0)$, $D(0)$, i.\,e. the terms independent of
% [sic] The German source calls these terms independent of n although the displayed value A(0)=n depends on n.
$n$, as they are appended.

Various relations among these functions are obtained in
the following way:
%CONT>>
